\chapter{Datos, validación y depuración}\label{cap:datos}

\section{Fuente y procedencia}

El profesor entregó el fichero \var{practica.sav} (SPSS), con \res{n-condados} condados y
\res{n-columnas} columnas; su huella SHA-256 empieza por \res{sha-corto}, lo que identifica sin
ambigüedad el fichero analizado. Procede del conjunto público \emph{OLS Regression Challenge}
\parencite{rippner2017}, que agrega datos de la \emph{American Community Survey}, de
\var{clinicaltrials.gov} y de \var{cancer.gov}. Antes de reutilizar la documentación pública de
las variables se comprobó, celda a celda y en el mismo orden de filas, que ambas fuentes
contienen los mismos datos. Sólo hay dos diferencias, y las dos están explicadas:
\begin{itemize}
  \item el fichero oficial ya trae como perdidos del sistema los casos anuales
  (\var{avgAnnCount}) de los condados cuyos registros no publican datos por condado, que el CSV
  público rellena con un valor centinela;
  \item el fichero oficial añade una variable sin documentar, \var{Notificadomuerte}, que resulta
  ser exactamente \var{avgAnnCount} $-$ \var{avgDeathsPerYear} en todos los condados en que
  ambas constan (regla R19).
\end{itemize}

La \cref{tab:variables} (\cref{ap:tablas}) recoge el diccionario: etiqueta, unidad, tipo, fuente
y papel de cada variable en el análisis.

\section{Variables que no pueden ser explicativas}

Algunas columnas se excluyen de los modelos por motivos de fondo, no por su significación
(\cref{tab:excluidas}, \cref{ap:tablas}):
\begin{itemize}
  \item \textbf{Fuga de información.} \var{avgDeathsPerYear} es el numerador de la tasa que se
  quiere explicar, y \var{Notificadomuerte} la contiene. Usarlas como explicativas daría un
  ajuste espectacular y vacío.
  \item \textbf{Escala.} \var{avgAnnCount} es un conteo que crece con la población; su
  versión por habitante es la incidencia, que sí se usa.
  \item \textbf{Redundancia exacta o casi exacta.} Las composiciones que suman \num{100}\,\%
  (educación de 18--24 años, raza) obligan a omitir una categoría de referencia, igual que se
  omite una indicadora al codificar una variable cualitativa; las edades medianas por sexo
  repiten la edad mediana total ($r > 0{,}9$).
\end{itemize}

\section{Validación}

Se definió un catálogo de \res{n-reglas} reglas que comprueban el dominio de cada variable
(porcentajes entre 0 y 100, tasas positivas), la coherencia entre variables (la edad mediana
total entre la masculina y la femenina; empleo más desempleo no mayor que 100\,\%), las
identidades contables (las cuatro categorías educativas de 18--24 años suman 100\,\%) y la
plausibilidad (valores repetidos con muchos decimales, mortalidad mayor que incidencia). La
\cref{tab:validacion} muestra el resultado antes y después de depurar: en el fichero original
fallan \res{n-reglas-error} reglas y \res{n-reglas-advertencia} emite una advertencia; tras la
depuración no queda ningún error (\res{n-reglas-error-dep}).

\begin{table}[htbp]
\centering
\small
\caption{Catálogo de reglas de validación: casos que las incumplen antes y después de la depuración}\label{tab:validacion}
\begin{tabularx}{\linewidth}{l X l r r}
\toprule
Id. & Regla & Severidad & Original & Depurado \\
\midrule
R01 & Identificador único & error & \num{0} & \num{0} \\
R02 & Porcentajes en [0, 100] & error & \num{0} & \num{0} \\
R03 & Tasas, renta y población positivas & error & \num{0} & \num{0} \\
R04 & Valores centinela & error & \num{206} & \num{0} \\
R05 & Edad mediana plausible ($\leq$ 100 años) & error & \num{30} & \num{0} \\
R06 & Coherencia de las edades medianas & error & \num{30} & \num{0} \\
R07 & Tamaño del hogar $\geq$ 1 persona & error & \num{61} & \num{0} \\
R08 & Identidad educativa de 18–24 años & error & \num{0} & \num{0} \\
R09 & Composición racial $\leq$ 100 \% & error & \num{0} & \num{0} \\
R10 & Empleo + desempleo $\leq$ 100 \% & error & \num{0} & \num{0} \\
R11 & Educación de 25+ años $\leq$ 100 \% & error & \num{0} & \num{0} \\
R12 & Jerarquía de coberturas sanitarias & error & \num{0} & \num{0} \\
R13 & Renta mediana dentro de su decil & error & \num{0} & \num{0} \\
R14 & Casos anuales $\geq$ muertes anuales & advert. & \num{0} & \num{0} \\
R15 & Razón mortalidad/incidencia plausible & advert. & \num{1} & \num{0} \\
R16 & Estado reconocido & error & \num{0} & \num{0} \\
R17 & Ausentes declarados & info. & \num{2504} & \num{872} \\
R18 & Inflación de ceros en ensayos clínicos & info. & \num{1931} & \num{1931} \\
R19 & Identidad de Notificadomuerte & error & \num{0} & \num{0} \\
\bottomrule
\end{tabularx}
\par\smallskip{\footnotesize\raggedright R17 y R18 son informativas: describen ausencias declaradas e inflación de ceros, que se tratan en etapas posteriores.\par}
\end{table}


\section{Depuración: once decisiones documentadas}

Cada incidencia se trata con una decisión registrada (\cref{tab:decisiones}) que expone el
problema, la evidencia que lo prueba con los propios datos, la acción y su justificación. Las
más relevantes son:

\begin{description}[style=nextline,leftmargin=1.2em]
  \item[Valor centinela en la incidencia (D03).] La incidencia vale exactamente
  \res{centinela-incidencia} en \res{n-centinela} condados, todos los de \res{estados-centinela}.
  Una medición con siete decimales no se repite por azar cientos de veces: es una imputación
  previa por la media disfrazada de dato. Se marca como ausente; conservarla introduciría cientos
  de observaciones sin variación en la incidencia y sesgaría su coeficiente hacia cero.
  \item[Edad mediana en meses (D04).] \res{n-edad-meses} condados declaran edades medianas de
  varios cientos de años. Divididas entre 12 quedan entre la mediana masculina y la femenina del
  mismo condado, lo que confirma el error de unidad; la razón media entre el valor erróneo y la
  media de las edades por sexo es \res{razon-edad} (desviación típica \res{dt-razon-edad}).
  \item[Tamaño del hogar dividido por cien (D05).] \res{n-hogar} condados declaran menos de una
  persona por hogar; multiplicados por 100 quedan en el rango del resto.
  \item[Recuperación por identidad (D06).] Falta el \res{pct-somecol}\,\% de
  \var{PctSomeCol18\_24}, pero las cuatro categorías educativas de 18--24 años suman 100\,\%
  (comprobado en los \res{n-identidad-completos} condados completos, con un error máximo de
  \res{error-identidad} puntos por redondeo), así que el valor ausente se calcula sin error.
  \item[Variable derivada (D07).] En \res{n-nativa} condados las cuatro categorías raciales
  publicadas suman menos del 80\,\%: falta la población nativa americana y multirracial. Se
  construye como residuo, \var{PctNativeMulti}.
  \item[Incidencia inverosímil (D08).] En \res{n-mir} condado (\res{condados-mir}) la
  mortalidad supera a la incidencia, cuando la razón mediana es \res{mir-mediana}; la incidencia
  de ese condado se marca como ausente.
  \item[Región censal (D09) y transformaciones (D10).] El estado tiene 51 categorías, algunas con
  uno o tres condados; se agrupa en las cuatro regiones de la Oficina del Censo. La población, la
  renta y los ensayos clínicos per cápita, muy asimétricos, se transforman con logaritmos.
\end{description}

\figura{depuracion}{Tres comprobaciones de la depuración: la edad mediana registrada en meses
queda, al dividirla entre 12, entre las medianas por sexo; el tamaño del hogar dividido por 100;
y la identidad educativa de 18--24 años que permite recuperar el dato ausente.}

{\footnotesize
\begin{xltabular}{\linewidth}{l p{0.28\linewidth} r X}
\caption{Registro de decisiones de depuración}\label{tab:decisiones} \\
\toprule
Id. & Decisión & Casos & Acción \\
\midrule
\endfirsthead
\multicolumn{4}{l}{\emph{(continuación)}} \\
\toprule
Id. & Decisión & Casos & Acción \\
\midrule
\endhead
\midrule
\multicolumn{4}{r}{\emph{(continúa)}} \\
\endfoot
\bottomrule

\endlastfoot
D01 & Geography separada en condado y estado & \num{3047} & Se separa por la última coma en Geography (condado) y state (estado). \\
D02 & Notificadomuerte excluida por fuga de información & \num{2841} & Se conserva en los datos pero se excluye de todos los modelos. \\
D03 & Valor centinela en la incidencia & \num{206} & Se sustituye por valor ausente. \\
D04 & Edad mediana registrada en meses & \num{30} & Se divide entre 12 y se verifica que el resultado quede entre las medianas masculina y femenina. \\
D05 & Tamaño medio del hogar dividido por 100 & \num{61} & Se multiplica por 100. \\
D06 & PctSomeCol18\_24 recuperada por identidad contable & \num{2285} & Se calcula como 100 \textminus{} (sin bachillerato + bachillerato + grado). \\
D07 & Población nativa o multirracial (variable derivada) & \num{86} & Se crea PctNativeMulti = 100 \textminus{} (blanca + negra + asiática + otra). \\
D08 & Incidencia inverosímil (mortalidad $\geq$ incidencia) & \num{1} & Se marca la incidencia como ausente. \\
D09 & Región censal & \num{3047} & Se agrupan en las cuatro regiones oficiales de la Oficina del Censo. \\
D10 & Transformaciones de variables muy asimétricas & \num{3047} & Población: logaritmo; renta: logaritmo; ensayos: log1p. \\
D11 & Variables excluidas de los modelos & \num{12} & Se excluyen: TARGET\_deathRate, avgDeathsPerYear, avgAnnCount, Notificadomuerte, binnedInc, Geography, state, MedianAgeMale, MedianAgeFemale, PctPrivateCoverageAlone, PctSomeCol18\_24, PctWhite. \\
\end{xltabular}
}


\section{Variables categóricas}\label{sec:categoricas}

El modelo incluye una explicativa cualitativa, la región censal (Sur, Noreste, Medio Oeste y
Oeste), que se codifica con tres variables indicadoras y el Sur como categoría de referencia
(la más numerosa: \res{n-sur} condados). Cada coeficiente de región es la diferencia esperada
de mortalidad respecto del Sur a igualdad del resto de variables, y su significación conjunta se
contrasta con un $F$ de tres grados de libertad. El estado no se usa como explicativa en el
modelo de efectos: con 51 categorías, varias casi vacías, consumiría 50 grados de libertad y
absorbería toda la variación entre estados; en cambio, se tiene en cuenta en la estimación de
los errores típicos (conglomerados) y como especificación de sensibilidad (efectos fijos de
estado, \cref{cap:sensibilidad}). El decil de renta (\var{binnedInc}) es una versión agrupada de
la renta mediana y se descarta en favor de la variable continua.
